% !TeX root = ../papers/nc-p3.tex
\section{二十個の二重点の極小分解から反対座標を復元する}
\label{sec:double-point-opposite-reconstruction}
$X\subset\mathbb P(U^*)$ は有限辺の第一射影とし、$S=\operatorname{Sym}U$ とする。
全階数の反対平方根から来る対称二次式行列を $M(x)$ と書く。
この呼称は十成分が $S_2$ の基底であり、$X$ 上で行列が射影的にベクトルの平方になることを表す。
$\mathbb P(S_2^*)$ の座標では $M$ の成分は独立な\emph{一次式}であり、この行列を $M(z)$ と書く。
\begin{theorem}\label{thm:double-point-opposite-reconstruction}
$M(x)$ の階数2以下の軌跡が有限であるとする。$X$ は非被約でもよい。
\begin{enumerate}
\item 三次小行列式は $\mathcal I_X^2$ の飽和イデアルを生成し、その極小分解は
\[
 0\to S(-10)^6\to S(-8)^{15}\to S(-6)^{10}\to I_{2X}\to0
\]
である。特に六次生成元と二次係数のシジジーは $X$ だけから決まる。
\item この分解の最左写像の各二次式を十個の二次座標 $z$ の一次式として読み直す。
得られる $15\times6$ 行列の階数3以下の\emph{台}は反対平方根の二次Veronese三次元多様体そのものである。
\item 同じ $X$ を通る別の全階数平方根 $N$ で $N(x)$ の一般階数が3以上なら、
そのVeronese像は先の $M$ のものに一致する。$N$ 側の行列式や階数2軌跡には仮定しない。
とくに $X$ が被約なら、元の座標と異なる反対平方根は単元倍を除いてただ一つである。
\end{enumerate}
\end{theorem}
\begin{proof}
任意の対称四行行列に対する次の複体を使う。添字を行列の型で示すと、
最初は交代行列 $D\mapsto DM$、次は跡零行列 $C\mapsto MC+C^{\mathsf t}M$、
最後は対称行列 $G\mapsto\operatorname{tr}(G\operatorname{adj}M)$ である。
二つの合成は零である。階数4または3の対称行列を合同対角化すれば三写像の階数は順に
$6,9,1$ である。例えば階数3では中央の像は最後の対角成分が零な対称行列全体になる。
階数降下は仮定によりアフィン錐で余次元3以上であり、
Buchsbaum--Eisenbudの複体の完全性判定が適用できる。
従って表示された次数の極小分解を得る。商は一次元Cohen--Macaulayなのでイデアルは飽和している。
Hilbert分子 $1-10t^6+15t^8-6t^{10}$ より次数は80である。
$X$ に沿って単元の一成分を消去すれば残る対称三行は $\mathcal I_X$ に成分を持ち、
三次小行列式は $\mathcal I_X^2$ に入る。$X$ は三次元滑らかな空間内の局所完全交叉なので
二重化の長さは $20+3\cdot20=80$。包含と次数から両イデアル層が一致する。

自由生成元の次数は各項で揃っているから極小分解は定数行列による基底変更を除き一意である。
最左写像は線型化後にも $D\mapsto D M(z)$ であり、その核次元は
$\binom{4-\operatorname{rank}M(z)}2$。従って指定された階数条件はちょうど
$\operatorname{rank}M(z)\le1$ を切り出す。

一意性では三次小行列式の\emph{線型シジジー}を十変数に戻すことが重要である。
$N(x)$ の各六次余因子も $I_{2X}$ に入り、その十生成元の定数線型結合となる。
この係数写像を $C_1$ とする。$N(x)$ に対する先の複体の中央の写像 $d_N$ は余因子の関係を与えるため、
完全な $M(x)$ の分解により $C_1 d_N=d_M C_2$ となる定数行列 $C_2$ がある。
これは各成分が二次式の等式であり、十変数の独立な一次式での等式へそのまま戻る。
十変数環 $P=\operatorname{Sym}(S_2)$ 上では $N(z),M(z)$ はそれぞれ一般対称行列の
座標変更である。それらの三次小行列式イデアル $J_N,J_M$ の関係は上と同じ複体で生成される
（対称行列の階数2以下は余次元3）。従って $C_1$ は次数零の $P$ 加群準同型
$J_N\to J_M$ を与える。$N(x)$ の少なくとも一つの余因子は非零なので $C_1\ne0$ で、
$J_M$ の三次十生成元は独立であるからこの準同型も零でない。
両イデアルの反射包は $P$ であり、次数零の準同型は反射包上のスカラー倍に尽きる。
これで $J_N\subseteq J_M$、同じ生成次数・生成元数により等号となる。
階数2以下の一般対称行列多様体の特異台は階数1以下である。
階数2では単元二行を消した残りが零という横断三式、階数1では対称三行の階数1以下の錐の頂点になる。
従って二つのVeronese像が一致する。被約の場合の残る一般階数2は
定理\ref{thm:opposite-root-triple-octics}の証明で排除され、一般階数1は元のVeronese像そのものである。
\end{proof}

\begin{corollary}\label{cor:normal-octics-general-finite-chart}
十五次元成分の稠密開集合で、各位置の二十点を三重に通る八次曲面は唯一で幾何学的に整かつ正規である。
各位置の二重化から上の最左写像で相手の座標配置が一意に復元される。
さらに任意の成分の被約有限窓でも、異なる反対根の八次式が正規な整曲面なら、
この復元と一意性が全ての特殊値に成立する。
\end{corollary}
\begin{proof}
階数2以下なら行列式の一次微分が全て零であるから、正規曲面ではこの軌跡は有限である。
一般点については既出の有理標本 $t$ の四位置で確認する。
\path{computations/verify_normal_configuration_octics.py} は既存の良い素数2657で二次評価空間を同定し、
各位置の二階偏導の $200\times165$ 評価の階数164を得る。
八次式は行列式から一つ実在し、標数零でも次元はちょうど1である。
明示平面に制限した行列式の三偏導に全十二次式を掛けると十九次全空間（次元210）を張る。
従ってこの平面切断は幾何学的に滑らかであり、曲面が可約、重複、または曲線に沿って特異となることはない。
超曲面のSerre判定で正規となる。これらの小行列条件と被約性は $t$ の開近傍で保たれる。
\end{proof}
